%%%%%%%%%%  scan idx 235  |  folio 222  %%%%%%%%%%
\section*{\underline{EXPOSÉ III}\\[3pt] \underline{\underline{CONDITIONS DE FINITUDE RELATIVES}}}
\begin{center}
\underline{par L. ILLUSIE}
\end{center}

\section*{1. Pseudo-cohérence relative}

\textbf{Théorème 1.1.} \uline{Soit} \(n\in\mathbb Z\). \uline{Pour chaque morphisme localement de type fini de schémas} \(f:X\longrightarrow Y\), \uline{il existe une et une seule sous-catégorie pleine} \(D(f)_{n\text{-}\coh}\) \uline{de} \(D(X)\)\SGAAuthorFootnote{1}{\(D(X)\) = catégorie dérivée de celle des \(\underline{O}_X\)-Modules.} \uline{telle que les conditions ci-dessous soient satisfaites} :
\begin{enumerate}[label=(\roman*)]
\item \uline{Si} \(E\in\operatorname{Ob} D(f)_{n\text{-}\coh}\), \uline{alors} \(i^*(E)\in\operatorname{Ob} D(fi)_{n\text{-}\coh}\) \uline{pour tout ouvert} \(i:U\to X\).

\item \uline{Si} \(E\in\ob D(X)\) \uline{et s'il existe un recouvrement ouvert} \((i_\alpha:U_\alpha\to X)_{\alpha\in A}\) \uline{tel que} \(i_\alpha^*(E)\in\ob D(fi_\alpha)_{n\text{-}\coh}\) \uline{pour tout} \(\alpha\in A\), \uline{alors} \(E\in\ob D(f)_{n\text{-}\coh}\).

\item \uline{Si}
\[
\begin{tikzcd}
X \arrow[rr,"f"] \arrow[dr,"h"'] & & Y \arrow[dl,"g"]\\
& Z &
\end{tikzcd}
\]
\uline{est un triangle commutatif de morphismes localement de type fini}, \uline{où} \(f\) \uline{est une immersion fermée et} \(g\) \uline{un morphisme lisse}, \uline{alors on a}
\[
  E\in\operatorname{Ob} D(h)_{n\text{-}\coh}
  \quad\Longleftrightarrow\quad
  f_*(E)\in D(Y)_{n\text{-}\coh}\qquad\mbox{\SGAAuthorFootMark{2}}\ .
\]
\SGAAuthorFootText{2}{\(D(Y)_{n\text{-}\coh}\) = sous-catégorie pleine de \(D(Y)\) formée des objets \(n\)-pseudo-cohérents (I~2.15).}
\end{enumerate}

\uline{Démonstration}. L'\uline{unicité} est évidente. Soit en effet \(f:X\to Y\) un morphisme localement de type fini. Choisissons un recouvrement ouvert \((i_\alpha:X_\alpha\to X)_{\alpha\in A}\) tel que, pour chaque \(\alpha\in A\), \(fi_\alpha\) se factorise en
\[
  X_\alpha\xrightarrow{i'_\alpha}Y_\alpha\xrightarrow{f'_\alpha}Y,
\]
où \(i'_\alpha\) est une immersion fermée et \(f'_\alpha\) un morphisme lisse. Il résulte aussitôt de (i), (ii), (iii) que les objets de \(D(f)_{n\text{-}\coh}\) sont les objets \(E\) de \(D(X)\) tels que
\[
  i'_{\alpha *}i_\alpha^*(E)\in \ob D(Y_\alpha)_{n\text{-}\coh};
\]
donc \(D(f)_{n\text{-}\coh}\) 
